% 67. hodina: Sbírka úloh. Úlohy 1–11 z fotografií uživatele, vlastní úlohy 12–17.
% Autorský přepis: české souřadnice [x;y], intervaly s českými závorkami.
\documentclass[11pt,a4paper]{article}
\usepackage[T1]{fontenc}\usepackage[utf8]{inputenc}\usepackage[czech]{babel}
\usepackage{lmodern,microtype,amsmath,amssymb,geometry,enumitem,tabularx,booktabs,tikz,multicol}
\usetikzlibrary{arrows.meta}
\geometry{top=16mm,bottom=17mm,left=18mm,right=18mm}
\pagestyle{plain}\setlength{\parindent}{0pt}\setlength{\parskip}{3pt}
\setlist[enumerate]{leftmargin=1.55em,topsep=2pt,itemsep=3pt}
\newcommand{\uloha}[2]{\par\needspace{4\baselineskip}\subsection*{#1\enspace #2}}
\usepackage{needspace}
\begin{document}
{\Large\bfseries Lineární funkce --- souhrnné procvičování}\hfill\textbf{67. hodina}
\par\vspace{1mm}Jméno a příjmení: \rule{62mm}{.4pt}\hfill Třída: \rule{19mm}{.4pt}
\smallskip
\textit{Sbírka není test a není nutné vyřešit všechny úlohy. Začni základními úlohami, poté si vyber obtížnější. Ke grafu vždy připiš předpis; bod zapisuj $[x;y]$. Můžeš využít přiložené souřadnicové sítě.}
\section*{A. Základní úlohy --- grafy a předpisy}
\uloha{1.}{Rychlé náčrty grafů lineárních funkcí}
Načrtněte grafy uvedených funkcí. Je-li graf přímkou, zvolte dva vhodné body. Uvažujte definiční obor $\mathbb R$.
\[
\begin{array}{r@{\quad}l@{\qquad\qquad}r@{\quad}l@{\qquad\qquad}r@{\quad}l}
\mathrm{a)}&y=1{,}5&\mathrm{b)}&y=3&\mathrm{c)}&y=-\pi\\[3pt]
\mathrm{d)}&y=2x&\mathrm{e)}&y=3x&\mathrm{f)}&y=-0{,}5x\\[3pt]
\mathrm{g)}&y=x+2&\mathrm{h)}&y=-2x+4&\mathrm{i)}&y=0{,}5x-1\\[3pt]
\mathrm{j)}&y=x-3&\mathrm{k)}&y=\frac32x-1&\mathrm{l)}&y=-\frac32x+1
\end{array}
\]
\textit{Navíc:} Rozhodněte, které grafy jsou vodorovné, které jsou rostoucí a které klesající.
\uloha{2.}{Předpis funkce z daných hodnot}
\begin{enumerate}[label=\alph*)]
\item Pro lineární funkci $h$ platí $h(3)=-5$ a $h(-1)=4$. Určete předpis $y=ax+b$ a načrtněte graf.
\item Je dáno $f(0)=1$ a $f(2)=5$. Najděte lineární funkci $f(x)=ax+b$ a načrtněte graf.
\item Lineární funkce $g$ splňuje $g(x+1)-g(x)=2$ pro každé $x\in\mathbb R$ a $g(0)=3$. Najděte předpis a načrtněte graf.
\end{enumerate}
\uloha{3.}{Předpis z bodů dvěma způsoby}
Lineární funkce prochází body $A[1;1]$ a $B[3;-2]$. Určete její předpis a sestrojte graf
\begin{enumerate}[label=\alph*)]
\item pomocí soustavy rovnic pro neznámé koeficienty $a,b$,
\item pomocí směrnice $a=\frac{\Delta y}{\Delta x}$ a absolutního členu $b$.
\end{enumerate}
\section*{B. Středně obtížné úlohy --- parametry a rodiny grafů}
\uloha{4.}{Směrnice z funkční hodnoty}
Pro $y=ax+b$ platí $b=-3$ a $f(2)=5$. Určete $a$, rozhodněte o monotonii a sestrojte graf.
\uloha{5.}{Předpis ze dvou funkčních hodnot}
O lineární funkci $f$ víme, že $f(1)=1{,}5$ a $f(-2)=-9$. Určete předpis $y=ax+b$ a rozhodněte, zda je rostoucí, či klesající.
\uloha{6.}{Rodina přímek se stejnou směrnicí}
Do jedné soustavy souřadnic zakreslete grafy
\[ y=0{,}7x+b,\qquad b\in\{-3;\,-1{,}5;\,0;\,1{,}5;\,2\}. \]
Co mají přímky společné a co se mění?
\uloha{7.}{Rodina přímek se stejným průsečíkem s osou $y$}
Do jedné soustavy souřadnic zakreslete grafy
\[ y=ax+3,\qquad a\in\{0;\,-3;\,-4;\,1{,}4;\,5\}. \]
Co mají grafy společné? Který je konstantní, které rostou a které klesají?
\section*{C. Náročné grafické úlohy --- po částech}
\uloha{8.}{Po částech lineární funkce z vlastností}
Nakreslete graf funkce $f$ se zadáním
\[
\begin{gathered}
D(f)=\langle-5;+\infty),\quad f(-5)=0,\quad f(0)=5,\quad f(4)=-5,\\
\text{$f$ na $\langle-5;0\rangle$ roste, na $\langle0;4\rangle$ klesá a na $(4;+\infty)$ je konstantní.}
\end{gathered}
\]
Graf vytvořte z úseček a vodorovné polopřímky. Zapište odpovídající předpis po částech a určete obor hodnot a maximum.
\uloha{9.}{Náčrty grafů po částech}
Načrtněte grafy a vyznačte správně krajní body jednotlivých částí:
\[
 f_1(x)=\begin{cases}x,&x\ge0,\\-x,&x<0,\end{cases}
\qquad
 f_2(x)=\begin{cases}1,&x\ge2,\\2x+1,&x<2.\end{cases}
\]
U obou určete $D(f)$ a $H(f)$. U $f_2$ rozhodněte, co se děje v $x=2$.
\section*{D. Další úlohy k výběru (vlastní doplnění)}
\uloha{10.}{Předpis z vyznačených grafů}
Na obrázku jsou tři přímky: plná $p$, čárkovaná $q$ a tečkovaná $r$. Určete jejich předpisy, směrnice a průsečíky s osou $y$.
\begin{center}
\begin{tikzpicture}[x=.74cm,y=.70cm,>=Latex]
\draw[step=1,gray!23,very thin] (-5,-4) grid (5,5);
\draw[->] (-5.35,0)--(5.45,0) node[right] {$x$};
\draw[->] (0,-4.35)--(0,5.45) node[above] {$y$};
\foreach \v in {-4,-3,-2,-1,1,2,3,4,5}{\draw (\v,.07)--(\v,-.07) node[below,font=\scriptsize] {$\v$};\draw (.07,\v)--(-.07,\v) node[left,font=\scriptsize] {$\v$};}
\begin{scope}\clip (-5,-4) rectangle (5,5);
\draw[very thick] (-5,6)--(5,-4);
\draw[very thick,dashed] (-5,-.5)--(5,4.5);
\draw[very thick,densely dotted] (-5,-2)--(5,-2);
\end{scope}
\node[fill=white,inner sep=1pt] at (3.5,-2.9) {$p$};
\node[fill=white,inner sep=1pt] at (2.1,3.5) {$q$};
\node[fill=white,inner sep=1pt] at (-3.7,-1.65) {$r$};
\end{tikzpicture}
\end{center}
\uloha{11.}{Parametr z intervalu}
Uvažujte rodinu funkcí $f_b(x)=2x+b$, kde $b\in\langle-2;1\rangle$.
Nakreslete obě krajní přímky, přidejte alespoň tři další a určete, kolik různých funkcí rodina obsahuje. Jak se mění průsečík s osou $y$?
\uloha{12.}{Omezený definiční obor}
Je dána funkce $g(x)=-2x+4$ pro $x\in\langle0;3)$. Sestrojte graf, označte krajní body a určete definiční obor, obor hodnot, nulový bod a průsečík s osou $y$.
\uloha{13.}{Sklon a rovnoběžnost}
Najděte předpis lineární funkce procházející bodem $[2;-1]$, jejíž graf je rovnoběžný s grafem $y=\frac32x+4$. Určete průsečíky s oběma osami.
\uloha{14.}{Dva tarify}
Tarif A má měsíční paušál 180 Kč a cenu 3 Kč za minutu. Tarif B má paušál 60 Kč a cenu 5 Kč za minutu. Pro počet minut $m\in\mathbb N_0$ sestavte funkce cen, porovnejte je graficky i početně a zjistěte, při jakém počtu minut stojí stejně.
\uloha{15.}{Vztah předpisu a jeho vlastností}
Pro $f(x)=ax+2$, $a\in\{-2;0;3\}$, napište všechny předpisy a u každého rozhodněte o monotonii, prostotě a oboru hodnot na $\mathbb R$.
\uloha{16.}{Nulový bod jako parametr}
Najděte všechna $b\in\mathbb R$, pro která má funkce $f_b(x)=2x+b$ nulový bod $x=3$. Popište graf této funkce.
\uloha{17.}{Je to vůbec lineární funkce na $\mathbb R$?}
Rozhodněte, zda funkce $u(x)=\frac{x^2-4}{x-2}$ je lineární na celém $\mathbb R$ a jak se její graf liší od přímky $y=x+2$.
\par\vspace{4mm}\hrule\vspace{2mm}
\textit{Při dokončení si ověř:} Jsou body správně ve tvaru $[x;y]$? Používám vhodné měřítko? Přímka je celá pouze tehdy, když to dovoluje definiční obor? Má více grafů v jedné soustavě jednoznačné popisky?
\end{document}
